\section{Bảng}
Bảng biểu được thể hiện như bảng~\ref{tab:my_label}, lưu ý flag \texttt{[H]} để disable floating (bảng được hiển thị đúng vị trí, không trôi lên đầu trang). Bảng~\ref{tab:my_label} là một trường hợp không sử dụng tag \texttt{[H]} và bảng bị trôi tít lên đầu trang:
\begin{table}[H]
	\centering
	\begin{tabular}{|l|l|}
		\hline
		\textbf{Tên con vật} & \textbf{Số chân} \\ \hline
		Gà                   & 2                \\ \hline
		Chó                  & 4                \\ \hline
		Trần Hoàng Tử        & 2                \\ \hline
	\end{tabular}
	\caption{Số chân của một số con vật, không có tag \texttt{[H]}}
	\label{tab:my_label}
\end{table}

Bảng~\ref{tab:my_label_with_H_tag} thể hiện bảng biểu với tag \texttt{[H]}\footnote{Tương tự cách sử dụng tag \texttt{[H]} với hình}. Để không phải mất thời gian tuổi trẻ ngồi chỉnh table, xài \href{https://www.tablesgenerator.com}{https://www.tablesgenerator.com}.

\begin{table}[H]
	\centering
	\begin{tabular}{|l|l|}
		\hline
		\textbf{Tên con vật} & \textbf{Số chân} \\ \hline
		Gà                   & 2                \\ \hline
		Chó                  & 4                \\ \hline
		Trần Hoàng Tử        & 2                \\ \hline
	\end{tabular}
	\caption{Số chân của một số con vật, có tag \texttt{[H]}}
	\label{tab:my_label_with_H_tag}
\end{table}.

Bảng~\ref{tab:my_label_with_tabularx} thể hiện bảng biểu với môi trường \texttt{tabularx}.

Môi trường này cho phép bạn định nghĩa các cột có độ rộng cố định hoặc co giãn, giúp bảng tự động xuống dòng khi nội dung quá dài. Bạn có thể sử dụng các ký hiệu \texttt{l} (left), \texttt{r} (right), \texttt{c} (center), \texttt{X} (co giãn) để định dạng các cột.

Có thể dùng \texttt{\textbackslash raggedleft}, \texttt{\textbackslash raggedright} và \texttt{\textbackslash centering} để canh lề nội dung trong các ô.

\begin{table}[H]
    \centering
    \begin{tabulary}{\textwidth}{|c|C|c|}
        \hline
        \textbf{STT} & \textbf{Nội dung}                                                 & \textbf{Số lượng} \\
        \hline
        1            & Đây là nội dung rất rất rất rất dài, LaTeX sẽ tự xuống dòng trong ô nhờ cột X & 10                \\
        \hline
        2            & Nội dung thứ 2, có thể dài nhiều dòng                              & 5                 \\
        \hline
    \end{tabulary}
    \caption{Tên bảng ở đây}
    \label{tab:my_label_with_tabularx}
\end{table}